\subsection{কৃত্রিম বুদ্ধিমত্তা}

\begin{learningbox}
এই অংশ শেষে শিক্ষার্থী পারবে:
\begin{itemize}
\item কৃত্রিম বুদ্ধিমত্তার ধারণা ব্যাখ্যা করতে।
\item AI, machine learning, expert system ও natural language processing-এর সম্পর্ক বুঝতে।
\item শিক্ষা, চিকিৎসা, কৃষি, ব্যবসা ও নিরাপত্তায় AI-এর ব্যবহার লিখতে।
\item AI ব্যবহারের সুবিধা, সীমাবদ্ধতা ও নৈতিক ঝুঁকি বিশ্লেষণ করতে।
\end{itemize}
\end{learningbox}

\subsubsection{কৃত্রিম বুদ্ধিমত্তার ধারণা}

মানুষ অভিজ্ঞতা থেকে শেখে, সমস্যা বুঝে সিদ্ধান্ত নেয় এবং ভাষা ব্যবহার করে যোগাযোগ করে। কম্পিউটারকে যখন এমনভাবে তৈরি করা হয় যে সে তথ্য বিশ্লেষণ করে শিখতে, যুক্তি প্রয়োগ করতে, সিদ্ধান্ত দিতে বা মানুষের ভাষা বুঝতে পারে, তখন সেটিকে কৃত্রিম বুদ্ধিমত্তা বলা হয়।

\begin{definitionbox}{সংজ্ঞা}
কৃত্রিম বুদ্ধিমত্তা বা Artificial Intelligence হলো কম্পিউটার বিজ্ঞানভিত্তিক এমন প্রযুক্তি, যার মাধ্যমে যন্ত্রকে মানুষের মতো শেখা, যুক্তি প্রয়োগ, সমস্যা সমাধান, ভাষা বোঝা, pattern শনাক্ত করা এবং সিদ্ধান্ত গ্রহণের ক্ষমতা দেওয়া হয়।
\end{definitionbox}

AI সবসময় মানুষের মতো চিন্তা করে না; বরং data, rule, algorithm ও model ব্যবহার করে সম্ভাব্য সঠিক ফল বের করে। যেমন, একটি spam filter হাজারো ই-মেইলের pattern দেখে কোনটি spam হতে পারে তা অনুমান করে।

\subsubsection{কৃত্রিম বুদ্ধিমত্তার বৈশিষ্ট্য}

\begin{keypointbox}{AI-এর গুরুত্বপূর্ণ বৈশিষ্ট্য}
\begin{itemize}
\item \textbf{Learning}: data থেকে pattern শেখে।
\item \textbf{Reasoning}: rule বা model ব্যবহার করে সিদ্ধান্ত দেয়।
\item \textbf{Problem solving}: জটিল সমস্যার সম্ভাব্য সমাধান খুঁজে।
\item \textbf{Perception}: image, sound, speech বা sensor data বুঝতে পারে।
\item \textbf{Language understanding}: মানুষের ভাষা বিশ্লেষণ ও response তৈরি করতে পারে।
\item \textbf{Adaptation}: নতুন data পেলে performance উন্নত করতে পারে।
\end{itemize}
\end{keypointbox}

\subsubsection{এক্সপার্ট সিস্টেম}

Expert system হলো AI-এর এমন একটি শাখা যেখানে কোনো বিশেষজ্ঞের জ্ঞান ও সিদ্ধান্তের নিয়ম computer program-এ সংরক্ষণ করা হয়। ব্যবহারকারী প্রশ্ন করলে system knowledge base ও inference engine ব্যবহার করে উত্তর বা পরামর্শ দেয়।

\begin{examplebox}{চিকিৎসা ক্ষেত্রে উদাহরণ}
একটি medical expert system রোগীর উপসর্গ, test report ও medical rule বিশ্লেষণ করে সম্ভাব্য রোগের তালিকা বা পরবর্তী test-এর পরামর্শ দিতে পারে। তবে এটি ডাক্তারের বিকল্প নয়; বরং decision support tool।
\end{examplebox}

\subsubsection{প্রাকৃতিক ভাষা প্রক্রিয়াকরণ}

Natural Language Processing বা NLP হলো AI-এর একটি অংশ, যা মানুষের ভাষা computer দিয়ে বোঝা, বিশ্লেষণ, অনুবাদ, summary তৈরি বা response দেওয়ার কাজ করে। Search engine, voice assistant, machine translation, chatbot, spelling correction এবং text summarization NLP-এর উদাহরণ।

\begin{mistakebox}{সাধারণ ভুল}
AI মানেই robot নয়। Robot একটি physical machine হতে পারে, কিন্তু AI একটি software-based decision বা learning technology। অনেক AI system কোনো robot ছাড়াই mobile app, website বা server-এ কাজ করে।
\end{mistakebox}

\subsubsection{মেশিন লার্নিং}

Machine Learning হলো AI-এর এমন পদ্ধতি যেখানে computer সরাসরি প্রতিটি rule লিখে না দিয়েও data থেকে pattern শেখে। উদাহরণস্বরূপ, অনেক ছবি দেখে model শিখতে পারে কোন ছবিতে handwritten digit আছে, অথবা কোন transaction fraud হতে পারে।

\begin{center}
\begin{tabular}{|p{0.30\textwidth}|p{0.56\textwidth}|}
\hline
\textbf{ধারণা} & \textbf{ব্যাখ্যা} \\
\hline
AI & যন্ত্রকে intelligent কাজ করানোর বিস্তৃত ক্ষেত্র \\
\hline
Machine Learning & data থেকে pattern শেখার পদ্ধতি \\
\hline
Deep Learning & neural network ব্যবহার করে জটিল pattern শেখা \\
\hline
NLP & মানুষের ভাষা বোঝা ও তৈরি করার প্রযুক্তি \\
\hline
\end{tabular}
\end{center}

\subsubsection{কৃত্রিম বুদ্ধিমত্তার ব্যবহারক্ষেত্র}

\begin{keypointbox}{ব্যবহারক্ষেত্র}
\begin{itemize}
\item \textbf{শিক্ষা}: adaptive learning, automated quiz checking, personalized feedback।
\item \textbf{চিকিৎসা}: disease prediction, medical image analysis, patient triage।
\item \textbf{কৃষি}: crop disease detection, weather-based advisory, smart irrigation।
\item \textbf{ব্যবসা}: recommendation system, fraud detection, chatbot।
\item \textbf{পরিবহন}: route optimization, traffic prediction, driver assistance।
\item \textbf{নিরাপত্তা}: face recognition, anomaly detection, cyber threat analysis।
\end{itemize}
\end{keypointbox}

\begin{examplebox}{বাংলাদেশ প্রেক্ষিত}
বাংলাদেশে কৃষি পরামর্শ app, mobile banking fraud detection, online customer support chatbot, e-learning platform এবং traffic monitoring system-এ AI-ভিত্তিক ধারণা ব্যবহার করা যেতে পারে বা আংশিকভাবে ব্যবহৃত হচ্ছে।
\end{examplebox}

\begin{keypointbox}{সুবিধা ও ঝুঁকি}
\begin{itemize}
\item সুবিধা: দ্রুত data analysis, repetitive কাজ automation, decision support, personalized service।
\item ঝুঁকি: data bias, privacy problem, ভুল সিদ্ধান্ত, job displacement, misinformation।
\item দায়িত্বশীল ব্যবহার: human supervision, transparent data use, privacy protection, fairness checking।
\end{itemize}
\end{keypointbox}

\begin{boardbox}{বোর্ড ফোকাস}
AI-এর প্রশ্নে সংজ্ঞা, বৈশিষ্ট্য, expert system, NLP, machine learning এবং ব্যবহারক্ষেত্র আলাদা করে লিখলে উত্তর পূর্ণ হয়। “AI মানেই robot”—এমন ভুল ধারণা এড়িয়ে চলবে।
\end{boardbox}

\begin{practicebox}
\textbf{MCQ}
\begin{enumerate}
\item Machine Learning কোন কাজের সঙ্গে সবচেয়ে বেশি সম্পর্কিত?\\
ক. Printer control \quad খ. data থেকে pattern শেখা \quad গ. শুধু typing \quad ঘ. monitor display
\item মানুষের ভাষা বোঝা ও বিশ্লেষণের AI শাখা কোনটি?\\
ক. NLP \quad খ. CAD \quad গ. GPS \quad ঘ. DBMS
\end{enumerate}

\textbf{সংক্ষিপ্ত প্রশ্ন}
\begin{enumerate}
\item Expert system কী?
\item AI ব্যবহারের দুটি নৈতিক ঝুঁকি লেখ।
\end{enumerate}
\end{practicebox}

\begin{sourcebox}
এই অংশে local HSC ICT PDF resources, Encyclopaedia Britannica-এর AI overview, IBM AI/ML explanation, Stanford AI Index-style ধারণা এবং online HSC ICT topic outlines মিলিয়ে নিজের ভাষায় লেখা হয়েছে।
\end{sourcebox}
